\documentclass[11pt]{article}
\usepackage[margin=1in]{geometry}
\usepackage{amsmath,amssymb,amsthm}
\usepackage[colorlinks=true,linkcolor=blue,urlcolor=blue,citecolor=blue]{hyperref}
\title{Geometric foundations for the eleven-square proof}
\author{}
\date{}
\begin{document}
\maketitle
\begin{abstract}
We give direct proofs of the nonavoidance and continuous-marker rules used
in the eleven-square reduction. All square interiors are open; center
regions are closed, and wall contact is allowed. The wall rule is proved
only in the parameter range needed for the reduction. These geometric
foundations are separate from the finite exclusions needed for optimality.
\end{abstract}

The first two rules correspond to Stromquist's Lemmas~1 and~2
\cite[ p.~2]{stromquist}. The triangle proof below also treats the boundary
and gives a finite-set generalization. The third rule is a directly proved variant of Lemma~4: it assumes
$b\ge1/2$ but permits every width $0<w<1$, whereas the published
lemma assumes $2\sqrt2-2<w<1$. Complete proofs are included, so the applications
depend on the stated hypotheses rather than an unchecked external theorem.

\section*{The two elementary nonavoidance rules}
\paragraph{Rectangle rule, with all boundaries included.}
Let $0<A,B\le1$. An open square of side $\lambda>1$ centered at
$(X,Y)\in[0,A]\times[0,B]$ either meets a coordinate axis or contains
$P=(A,B)$. Suppose it avoids both axes. Choose edge axes
$u=(c,s),v=(-s,c)$ with $c,s\ge0$, $c^2+s^2=1$, and put $\rho=c+s$.
Avoidance of both axes implies $X,Y\ge\lambda\rho/2$. Thus the
coordinates of $d=P-(X,Y)$ are nonnegative and at most
$M=1-\lambda\rho/2$. In a feasible instance $M\ge0$, and both edge
projections satisfy
\[
 |u\cdot d|,\ |v\cdot d|\le\rho M<\lambda/2.
\]
Indeed, the strict inequality is equivalent to
$2\rho<\lambda(\rho^2+1)$, which follows from
$\rho^2+1\ge2\rho>0$ and $\lambda>1$. Therefore $P$ belongs to the
open square, even if the square touches an axis or its center lies on
an edge of the center rectangle.

\paragraph{Triangle rule, including edges and vertices.}
Let a triangle have all side lengths at most one. An open square of
side $\lambda>1$ whose center lies in the \emph{closed} triangle
contains one of its vertices. In fact, the statement holds with a finite
set $V$ of diameter at most one in place of the three vertices and its
convex hull in place of the triangle.
Translate the square's center to zero and use its edge coordinates.
Suppose every point of $V$ avoids its open interior. Every point then
satisfies one of $x\ge\lambda/2$, $x\le-\lambda/2$,
$y\ge\lambda/2$, $y\le-\lambda/2$. Both opposite inequalities in either
coordinate cannot occur at different points, as they would give a pair
at distance at least $\lambda>1$. Reflecting coordinates if necessary,
all violations are therefore of types $x\ge\lambda/2$ and
$y\ge\lambda/2$.
If only one type occurs, that coordinate is strictly positive throughout
$V$, excluding zero from its convex hull. If both types occur, the
pairwise-distance bound forces every point to satisfy
$x> -\lambda/2$ and $y> -\lambda/2$. Each point thus satisfies $x+y>0$,
again excluding zero from the convex hull. This contradiction proves the
open-ownership conclusion without any separate center-edge argument.

\section*{A direct proof of the required wall-quadrilateral rule}
Let $0<w<1$, $1/2\le b<1$, and
\[
 w^2+(1-b)^2\le1.
\]
Suppose that for every $0<\theta<\pi/2$,
\begin{equation}\label{eq:independent-wall-hypothesis}
 b\le \frac{\cos\theta}{1+\cos\theta}
       +\frac{1-w\cos\theta}{\sin\theta}.
\end{equation}
Write
\[
 P=(0,1),\qquad R=(w,b),\qquad
 D=\operatorname{conv}\{(0,0),(w,0),R,P\}.
\]
Then the open interior of every square of side $\lambda>1$ whose center
lies in the closed quadrilateral $D$ meets the $x$-axis or contains $P$ or
$R$. This is a directly proved variant of Stromquist\textup{~\cite[Lemma~4, pp.~3--4]{stromquist}} in the range $b\ge1/2$; its wider width range is justified by the proof below. Its scalar hypothesis is stated for every angle, so no assertion about the location or uniqueness of a minimum is required.

\paragraph{Coordinates and boundary convention.}
Suppose, for a contradiction, that the square's open interior avoids all
three targets. Let its center be $Z=(X,Y)$. Choose its orthonormal edge
axes as
\[
 u=(c,s),\qquad v=(-s,c),\qquad c,s\ge0,\quad c^2+s^2=1,
\]
and put $\rho=c+s$. Such a choice covers every square orientation modulo
$\pi/2$. The center hypotheses give
\[
 0\le X\le w,\qquad 0\le Y\le1-(1-b)X/w\le1.
\]
The square's vertical half-width is $\lambda\rho/2$. Since it avoids the
$x$-axis and its center is above it,
\begin{equation}\label{eq:wall-bottom}
 Y\ge\lambda\rho/2.
\end{equation}
Equality, corresponding to a vertex or edge on the wall, is allowed.
Put $U=u\cdot Z$, $V=v\cdot Z$, and let the edge coordinates of the two
markers relative to $Z$ be
\[
 p=u\cdot(P-Z)=s-U,\quad q=v\cdot(P-Z)=c-V,
\]
\[
 r=u\cdot(R-Z)=cw+sb-U,\quad t=v\cdot(R-Z)=-sw+cb-V.
\]
Avoidance of an \emph{open} square means at least one coordinate has
absolute value at least $\lambda/2$. No strict outside assumption is made.

\paragraph{Only four face combinations remain.}
First, $q=c(1-Y)+sX\ge0$, whereas
\[
 p=s(1-Y)-cX\le s(1-\lambda\rho/2)<\lambda/2.
\]
The strict inequality follows from
$1+s\rho-2s=(1-s)^2+sc\ge0$ and $\lambda>1$.
Consequently the first marker can be avoided only through
\[
 \text{(A)}\ p\le-\lambda/2,
 \qquad\hbox{or}\qquad
 \text{(B)}\ q\ge\lambda/2.
\]
For the second marker,
\[
 r=c(w-X)+s(b-Y)\ge-s(1-b)\ge-1/2>-\lambda/2,
\]
using $b\ge1/2$, and
\[
 t=-s(w-X)+c(b-Y)\le c(1-\lambda\rho/2)<\lambda/2.
\]
Here $1+c\rho-2c=(1-c)^2+cs\ge0$ gives the final inequality.
Thus its only avoidance possibilities are
\[
 \text{(C)}\ r\ge\lambda/2,
 \qquad\hbox{or}\qquad
 \text{(D)}\ t\le-\lambda/2.
\]
If a marker violates two faces, choose either one; the four combinations
below exhaust all possibilities, including marker-on-boundary cases.

\paragraph{Three combinations are impossible.}
Set $\delta=1-b$ and
\[
 d=u\cdot(R-P)=cw-s\delta,
 \qquad E=-v\cdot(R-P)=sw+c\delta>0.
\]
Then $r=p+d$, $t=q-E$, and
\begin{equation}\label{eq:wall-short-top}
 d^2+E^2=w^2+\delta^2\le1<\lambda^2.
\end{equation}
Combination (A),(C) gives $d=r-p\ge\lambda$, contradicting
\eqref{eq:wall-short-top}. Combination (B),(D) similarly gives
$E=q-t\ge\lambda$.

For combination (A),(D), (A) yields $U\ge s+\lambda/2$.
Since the maximum of $u\cdot Z$ on $D$ is $\max(s,cw+sb)$,
we get $d\ge\lambda/2>0$. In particular $0<d,E<\lambda$ by
\eqref{eq:wall-short-top}. The top-edge condition on the center says
\[
 (\delta,w)\cdot(P-Z)\ge0.
\]
Since $(\delta,w)=E u+d v$, this is $Ep+dq\ge0$.
But (A),(D) give $p\le-\lambda/2$ and $q\le E-\lambda/2$, so
\[
 Ep+dq\le dE-\frac\lambda2(d+E)<0.
\]
The last inequality uses $d,E<\lambda$:
$2dE<\lambda E+\lambda d$. This is a contradiction.
This step is exactly where the center's position below the sloping edge
is needed; it cannot be omitted or replaced by a coarse rectangle.

\paragraph{The remaining combination gives the scalar condition.}
In combination (B),(C),
\[
 V\le c-\lambda/2,\qquad U\le cw+sb-\lambda/2.
\]
Since $Y=sU+cV$, these inequalities together with
\eqref{eq:wall-bottom} imply
\begin{equation}\label{eq:wall-final-inequality}
 \lambda(c+s)\le c^2+scw+s^2b.
\end{equation}
If $s=0$, this is $\lambda\le1$, impossible. If $c=0$, it is
$\lambda\le b<1$, also impossible. Otherwise $s,c>0$ and
$\lambda>1$ make \eqref{eq:wall-final-inequality} imply
\[
 b>\frac{c+s-c^2-scw}{s^2}
   =\frac{c}{1+c}+\frac{1-wc}{s},
\]
contradicting \eqref{eq:independent-wall-hypothesis}.
Thus the open square meets the wall or owns an endpoint, as claimed.
All center boundaries, marker boundaries, wall contacts, and the two
axis-aligned orientations have been included explicitly.


\paragraph{The initial wall parameters.}
For the twelve-marker cover, take $w=47/50$ and $b=3/4$.
The short top-edge condition is
$(47/50)^2+(1/4)^2=9461/10000<1$.
With $z=\tan(\theta/2)\in(0,1)$, the scalar hypothesis follows from
\[
100z\left(\frac{\cos\theta}{1+\cos\theta}
 +\frac{1-w\cos\theta}{\sin\theta}-b\right)
=3-25z+97z^2-50z^3>0.
\]
For $z\le1/2$ the polynomial is at least $3-25z+72z^2$, whose
discriminant is $-239$. For $z\ge1/2$ it is at least
$3-25z+47z^2$, which is increasing there and equals $9/4$ at $1/2$.
Other marker positions require their own scalar checks.
If the left endpoint has height $H$ instead of one, scale all lengths by
$H$; the required strict side inequality is then $a/H>1$.

\section*{Continuous markers preserve ownership}
The following generalized form of Bentz's Theorem~8\cite[pp.~3--4]{bentz} applies to
fixed open sets, and in particular to the interiors of packed squares.

Let $B_1,\ldots,B_m$ be pairwise disjoint open sets, and let
$p_1,\ldots,p_n:[0,1]\to\mathbb R^2$ be continuous curves.
Suppose every $B_i$ contains at least one $p_j(t)$ at every time $t$.
Keep each initially unowned marker fixed, and also keep every marker
fixed whose initial owner contains another marker. Then every initially
owned marker stays in its initial owner for the entire interval.

To prove this, suppose an exit occurs. For each marker initially in an
owner, its set of exit times is closed; hence, among the finitely many
markers that exit, there is a first exit time $\tau>0$. Let $p_j$ leave
$B_i$ then. At time $\tau$, another marker $p_k$ lies strictly inside
$B_i$, because that set must still contain a marker. Openness and
continuity put $p_k(t)$ in $B_i$ just before $\tau$ as well.
This marker cannot have been initially unowned: it would have been fixed
outside every owner. If it was initially in $B_l$, minimality of $\tau$
puts it in $B_l$ just before $\tau$. Disjointness forces $l=i$.
Thus $p_j$ and $p_k$ initially shared an owner, so $p_j$ was fixed and
could not exit. This is a contradiction. Simultaneous exits and points
on square boundaries are covered by the same argument.

For a continuously unavoidable marker configuration, the premise about
every fixed packed square holds automatically. If an owner is convex,
it consequently contains the convex hull of the entire path of any
marker it owns. The moving objects in this argument are the markers;
the packed squares stay fixed. In particular, the freezing assumptions
must be checked before a marker deformation can force a new owned region.

\begin{thebibliography}{2}
\bibitem{stromquist}
W.~Stromquist, \emph{Packing 10 or 11 Unit Squares in a Square},
Electronic Journal of Combinatorics \textbf{10} (2003), R8,
\href{https://doi.org/10.37236/1701}{doi:10.37236/1701}.
The relevant statements are Lemmas~1, 2 and~4; no other result from that
paper is needed for the proofs above.
\bibitem{bentz}
W.~Bentz, \emph{Optimal Packings of 22 and 33 Unit Squares in a Square},
\href{https://arxiv.org/abs/1606.03746v1}{arXiv:1606.03746v1}, Theorem~8,
pp.~3--4. The versioned record was submitted in 2016; its PDF title page
is dated 20 October 2018.
\end{thebibliography}
\end{document}
